\chapter{COMBINAÇÕES UTILIZADAS NA ESTIMATIVA DA PROFUNDIDADE POR HVSR}\label{ap:hvsr}
\setcounter{figure}{0}\renewcommand{\thefigure}{F.\arabic{figure}}
\setcounter{table}{0}\renewcommand{\thetable}{F.\arabic{table}}

Este apêndice apresenta as 18 combinações utilizadas na estimativa da profundidade associada ao pico fundamental em cada data. Foram combinados os limites inferior, central e superior da frequência fundamental, dois valores de velocidade de onda P e três coeficientes de Poisson. As linhas correspondentes à frequência central, definida pela mediana lognormal, são destacadas para facilitar sua identificação (Tabelas F.1 e F.2).

{\scriptsize
\setlength{\tabcolsep}{2pt}
\renewcommand{\arraystretch}{1.08}
\begin{longtable}{@{} >{\raggedright\arraybackslash}p{0.200\linewidth} >{\centering\arraybackslash}p{0.140\linewidth} >{\centering\arraybackslash}p{0.140\linewidth} >{\centering\arraybackslash}p{0.140\linewidth} >{\centering\arraybackslash}p{0.140\linewidth} >{\centering\arraybackslash}p{0.140\linewidth} @{}}
  \caption{Cenários de profundidade por HVSR para 01/02/2020.}\label{tab:resultado-F-1} \\
    \toprule
    \textbf{Dia} & \textbf{Frequência (Hz)} & \textbf{Velocidade P (m/s)} & \textbf{Coeficiente de Poisson (ν)} & \textbf{Velocidade S (m/s)} & \textbf{Profundidade HVSR (m)} \\
    \midrule
    \endfirsthead
  \caption[]{Cenários de profundidade por HVSR para 01/02/2020. (continuação)} \\
    \toprule
    \textbf{Dia} & \textbf{Frequência (Hz)} & \textbf{Velocidade P (m/s)} & \textbf{Coeficiente de Poisson (ν)} & \textbf{Velocidade S (m/s)} & \textbf{Profundidade HVSR (m)} \\
    \midrule
    \endhead
    \midrule \multicolumn{6}{r}{\textit{continua}} \\
    \endfoot
    \bottomrule
    \endlastfoot
    01/02/2020 & f₋₁σ = 0,813 & 1.617,06 & 0,46 & 440,11 & 135,37 \\
    \midrule
    01/02/2020 & f₋₁σ = 0,813 & 1.617,06 & 0,47 & 384,72 & 118,34 \\
    \midrule
    01/02/2020 & f₋₁σ = 0,813 & 1.617,06 & 0,48 & 317,13 & 97,55 \\
    \midrule
    01/02/2020 & f₋₁σ = 0,813 & 1.649,58 & 0,46 & 448,96 & 138,10 \\
    \midrule
    01/02/2020 & f₋₁σ = 0,813 & 1.649,58 & 0,47 & 392,46 & 120,72 \\
    \midrule
    01/02/2020 & f₋₁σ = 0,813 & 1.649,58 & 0,48 & 323,51 & 99,51 \\
    \midrule
    01/02/2020 & f₀ = 0,857 & 1.617,06 & 0,46 & 440,11 & 128,39 \\
    \midrule
    01/02/2020 & f₀ = 0,857 & 1.617,06 & 0,47 & 384,72 & 112,23 \\
    \midrule
    01/02/2020 & f₀ = 0,857 & 1.617,06 & 0,48 & 317,13 & 92,51 \\
    \midrule
    01/02/2020 & f₀ = 0,857 & 1.649,58 & 0,46 & 448,96 & 130,97 \\
    \midrule
    01/02/2020 & f₀ = 0,857 & 1.649,58 & 0,47 & 392,46 & 114,49 \\
    \midrule
    01/02/2020 & f₀ = 0,857 & 1.649,58 & 0,48 & 323,51 & 94,37 \\
    \midrule
    01/02/2020 & f₊₁σ = 0,904 & 1.617,06 & 0,46 & 440,11 & 121,76 \\
    \midrule
    01/02/2020 & f₊₁σ = 0,904 & 1.617,06 & 0,47 & 384,72 & 106,44 \\
    \midrule
    01/02/2020 & f₊₁σ = 0,904 & 1.617,06 & 0,48 & 317,13 & 87,74 \\
    \midrule
    01/02/2020 & f₊₁σ = 0,904 & 1.649,58 & 0,46 & 448,96 & 124,21 \\
    \midrule
    01/02/2020 & f₊₁σ = 0,904 & 1.649,58 & 0,47 & 392,46 & 108,58 \\
    \midrule
    01/02/2020 & f₊₁σ = 0,904 & 1.649,58 & 0,48 & 323,51 & 89,50 \\
\end{longtable}
}

{\scriptsize
\setlength{\tabcolsep}{2pt}
\renewcommand{\arraystretch}{1.08}
\begin{longtable}{@{} >{\raggedright\arraybackslash}p{0.200\linewidth} >{\centering\arraybackslash}p{0.140\linewidth} >{\centering\arraybackslash}p{0.140\linewidth} >{\centering\arraybackslash}p{0.140\linewidth} >{\centering\arraybackslash}p{0.140\linewidth} >{\centering\arraybackslash}p{0.140\linewidth} @{}}
  \caption{Cenários de profundidade por HVSR para 29/02/2020.}\label{tab:resultado-F-2} \\
    \toprule
    \textbf{Dia} & \textbf{Frequência (Hz)} & \textbf{Velocidade P (m/s)} & \textbf{Coeficiente de Poisson (ν)} & \textbf{Velocidade S (m/s)} & \textbf{Profundidade HVSR (m)} \\
    \midrule
    \endfirsthead
  \caption[]{Cenários de profundidade por HVSR para 29/02/2020. (continuação)} \\
    \toprule
    \textbf{Dia} & \textbf{Frequência (Hz)} & \textbf{Velocidade P (m/s)} & \textbf{Coeficiente de Poisson (ν)} & \textbf{Velocidade S (m/s)} & \textbf{Profundidade HVSR (m)} \\
    \midrule
    \endhead
    \midrule \multicolumn{6}{r}{\textit{continua}} \\
    \endfoot
    \bottomrule
    \endlastfoot
    29/02/2020 & f₋₁σ = 0,852 & 1.617,06 & 0,46 & 440,11 & 129,10 \\
    \midrule
    29/02/2020 & f₋₁σ = 0,852 & 1.617,06 & 0,47 & 384,72 & 112,85 \\
    \midrule
    29/02/2020 & f₋₁σ = 0,852 & 1.617,06 & 0,48 & 317,13 & 93,02 \\
    \midrule
    29/02/2020 & f₋₁σ = 0,852 & 1.649,58 & 0,46 & 448,96 & 131,69 \\
    \midrule
    29/02/2020 & f₋₁σ = 0,852 & 1.649,58 & 0,47 & 392,46 & 115,12 \\
    \midrule
    29/02/2020 & f₋₁σ = 0,852 & 1.649,58 & 0,48 & 323,51 & 94,89 \\
    \midrule
    29/02/2020 & f₀ = 0,873 & 1.617,06 & 0,46 & 440,11 & 126,03 \\
    \midrule
    29/02/2020 & f₀ = 0,873 & 1.617,06 & 0,47 & 384,72 & 110,17 \\
    \midrule
    29/02/2020 & f₀ = 0,873 & 1.617,06 & 0,48 & 317,13 & 90,82 \\
    \midrule
    29/02/2020 & f₀ = 0,873 & 1.649,58 & 0,46 & 448,96 & 128,57 \\
    \midrule
    29/02/2020 & f₀ = 0,873 & 1.649,58 & 0,47 & 392,46 & 112,39 \\
    \midrule
    29/02/2020 & f₀ = 0,873 & 1.649,58 & 0,48 & 323,51 & 92,64 \\
    \midrule
    29/02/2020 & f₊₁σ = 0,894 & 1.617,06 & 0,46 & 440,11 & 123,04 \\
    \midrule
    29/02/2020 & f₊₁σ = 0,894 & 1.617,06 & 0,47 & 384,72 & 107,56 \\
    \midrule
    29/02/2020 & f₊₁σ = 0,894 & 1.617,06 & 0,48 & 317,13 & 88,66 \\
    \midrule
    29/02/2020 & f₊₁σ = 0,894 & 1.649,58 & 0,46 & 448,96 & 125,52 \\
    \midrule
    29/02/2020 & f₊₁σ = 0,894 & 1.649,58 & 0,47 & 392,46 & 109,72 \\
    \midrule
    29/02/2020 & f₊₁σ = 0,894 & 1.649,58 & 0,48 & 323,51 & 90,45 \\
\end{longtable}
}
